\section{Background}\label{sec:background}

\subsection{Chips, Slices, and Hosts}\label{sec:bg-tpu}
A TPU v5e chip has 16~GiB of high-bandwidth memory (HBM)~\cite{googletpuv5e}. % src: googletpuv5e as read on 2026-11-23 ("16 GiB"), paper-v3/SOURCE_NOTES.md section 1
The chips of a slice are linked by an inter-chip interconnect in a two-dimensional torus~\cite{googletpuv5e}. Our
slice had eight chips in a 2~$\times$~4 topology. % src: configs/af_spike_job.yaml -> nodeSelector; recomputed: V2.topology
A host is the CPU machine that runs the Python process and drives the chips. JAX treats each chip as one device of that
process. In this paper, a chip is one TPU v5e chip, and a slice is the set of eight chips that each Kubernetes Job received.

\subsection{Compilation and Mapping in JAX}\label{sec:bg-jax}
Before it compiles a function, JAX traces it into an intermediate program. \texttt{jax.jit} compiles a function with
XLA when it is first called, caches the compiled code, and reuses it in later calls~\cite{jaxjit}. In this paper,
compilation means tracing and XLA compilation together. We do not separate the two. JAX can also store compiled
programs on disk, so that a later run can load them instead of compiling again~\cite{jaxcache}. We call this store the
compilation cache.

JAX offers several ways to run more than one query. \texttt{jax.vmap} maps a function over a batch axis of its
inputs on one device~\cite{jaxvmap}. \texttt{jax.pmap} replicates a function and runs each replica on its own device. Each device
receives one element of the leading axis of the input~\cite{jaxpmap}. The documentation of \texttt{jax.pmap} now
recommends \texttt{jax.shard\_map} instead~\cite{jaxpmap}. \texttt{jax.lax.pmean} computes the mean of an array over the
devices of a mapped axis~\cite{jaxpmean}.

JAX can also partition one computation automatically over a mesh, a grid of devices with named axes. The GSPMD partitioner of XLA starts from
the sharding annotations that the program places on some arrays. It derives the placement of the other arrays from
them~\cite{xu2021gspmd}. The OpenXLA guide to Shardy for JAX users says that JAX now always uses the Shardy
partitioner~\cite{shardy2026}.
We read this guide on 23 November 2026. We did not check which partitioner JAX 0.10.2 used. % src: shardy2026 ("read 2026-11-23", "Shardy is now always enabled in JAX"), paper-v3/SOURCE_NOTES.md section 2; recomputed: SEP.cpu.0915.jax
In Shardy, an array is replicated along each mesh axis that does not shard one of its dimensions~\cite{shardyrepr}.

\subsection{AlphaFold2}\label{sec:bg-af2}
AlphaFold2 reads a multiple sequence alignment (MSA), which lines up the query with related protein sequences. A stack
of Evoformer blocks updates a representation of the MSA and a representation of residue pairs. Residues are the amino
acids of the protein chain. A structure module then turns these representations into atom
coordinates~\cite{jumper2021alphafold}. The public implementation uses Haiku, a neural-network library for
JAX~\cite{haiku2020github}. The
model can recycle its output as a new input for a fixed number of
iterations. The public implementation runs recycling in a \texttt{while\_loop}.
% One more line in this paragraph, so that the sentence "Table I does not show them" (Section III-A) is not split
% by Table I. The text does not change.
\looseness=1

The implementation can also run an ensemble. Its members are copies of the same input, featurized differently.
Featurization turns the sequence and the MSA into the arrays that the model reads. The implementation computes the
first member outside a loop and the others in a sequential Haiku \texttt{while\_loop}~\cite{alphafoldrepo}. % src: AlphaFold2 at commit c77e5d2, alphafold/model/modules.py:174-205 ("hk.while_loop"), paper-v3/SOURCE_NOTES.md section 3; recomputed: SEP.cpu.0915.alphafold_commit
The authors of AlphaFold2 turn off ensembling for most inference~\cite{jumper2021alphafold}. In the released code, the
preset \texttt{monomer\_casp14} uses eight ensemble members, and the other presets use one~\cite{alphafoldrepo}. % src: AlphaFold2 at commit c77e5d2, run_alphafold.py:601-604 ("num_ensemble = 8", "num_ensemble = 1"); jumper2021alphafold ("we turn off ensembling for most inference")
Our runs used one ensemble member, so this loop ran no iteration in the timed calls. % src: src/spike_tpu_forward_pass.py ("cfg.data.eval.num_ensemble = 1")
The \texttt{predict} method of the model runs \texttt{apply}, the forward pass on the device. It then computes
confidence metrics on the host with NumPy.
